\textbf{Phần 1. Trắc nghiệm lựa chọn}
\setcounterpageref{ex}{0}

\begin{ex} %[5H3V16]
Một căn phòng hình chữ nhật có chiều dài $6\text{ m}$ và chiều rộng $4{,}5\text{ m}$. Người ta dùng các viên gạch men hình vuông cạnh $30\text{ cm}$ để lát sàn. Biết giá tiền mỗi viên gạch là $25\,000$ đồng. Tổng số tiền mua gạch để lát kín căn phòng là bao nhiêu triệu đồng?
\choice
{$6{,}5$ triệu đồng}
{\True $7{,}5$ triệu đồng}
{$8$ triệu đồng}
{$7$ triệu đồng}
\loigiai{
Diện tích căn phòng là:
\[
6 \times 4{,}5 = 27\text{ (m}^2\text{)} = 270\,000\text{ cm}^2.
\]
Diện tích một viên gạch men là:
\[
30 \times 30 = 900\text{ (cm}^2\text{)}.
\]
Số viên gạch cần dùng là:
\[
270\,000 : 900 = 300\text{ (viên)}.
\]
Tổng số tiền mua gạch là:
\[
300 \times 25\,000 = 7\,500\,000\text{ (đồng)} = 7{,}5\text{ triệu đồng}.
\]
\textbf{Đáp án của câu hỏi là B}
}
\end{ex}

\begin{ex} %[5H3V16]
Một tấm bìa hình vuông có diện tích là $1\text{ m}^2$. Người ta cắt tấm bìa đó thành các hình vuông nhỏ có cạnh $1\text{ dm}$. Sau đó đem các hình vuông nhỏ này xếp nối tiếp nhau thành một hàng thẳng. Chiều dài của hàng thẳng đó là
\choice
{$1\text{ m}$}
{\True $10\text{ m}$}
{$100\text{ m}$}
{$10\text{ dm}$}
\loigiai{
Vì $1\text{ m}^2 = 100\text{ dm}^2$ nên diện tích tấm bìa ban đầu gồm $100$ hình vuông nhỏ cạnh $1\text{ dm}$.\\
Khi xếp nối tiếp $100$ hình vuông nhỏ cạnh $1\text{ dm}$ thành một hàng thẳng, chiều dài hàng thẳng là:
\[
100 \times 1\text{ dm} = 100\text{ dm} = 10\text{ m}.
\]
\textbf{Đáp án của câu hỏi là B}
}
\end{ex}

\begin{ex} %[5H3V16]
Số thích hợp điền vào chỗ chấm để $5\text{ m}^2\ 8\text{ cm}^2 = \dots\text{ cm}^2$ là
\choice
{$580$}
{$5\,080$}
{\True $50\,008$}
{$500\,008$}
\loigiai{
Vì $1\text{ m}^2 = 10\,000\text{ cm}^2$ nên:
\[
5\text{ m}^2\ 8\text{ cm}^2 = 50\,000\text{ cm}^2 + 8\text{ cm}^2 = 50\,008\text{ cm}^2.
\]
\textbf{Đáp án của câu hỏi là C}
}
\end{ex}

\begin{ex} %[5H3V16]
Một bức tường hình chữ nhật có kích thước $4\text{ m} \times 3\text{ m}$. Người ta lắp một cửa sổ hình vuông cạnh $12\text{ dm}$ trên bức tường đó. Diện tích phần tường cần sơn (trừ diện tích cửa sổ) là
\choice
{$10{,}56\text{ m}^2$}
{\True $10{,}56\text{ m}^2$}
{$11{,}88\text{ m}^2$}
{$9{,}56\text{ m}^2$}
\loigiai{
Diện tích cả bức tường là:
\[
4 \times 3 = 12\text{ (m}^2\text{)}.
\]
Đổi cạnh cửa sổ: $12\text{ dm} = 1{,}2\text{ m}$.\\
Diện tích cửa sổ hình vuông là:
\[
1{,}2 \times 1{,}2 = 1{,}44\text{ (m}^2\text{)}.
\]
Diện tích phần tường cần sơn là:
\[
12 - 1{,}44 = 10{,}56\text{ (m}^2\text{)}.
\]
\textbf{Đáp án của câu hỏi là A}
}
\end{ex}

\begin{ex} %[5H3V16]
Một tấm tôn hình thoi có độ dài hai đường chéo là $8\text{ dm}$ và $65\text{ cm}$. Diện tích của tấm tôn đó tính theo đơn vị đề-xi-mét vuông là
\choice
{$52\text{ dm}^2$}
{\True $26\text{ dm}^2$}
{$260\text{ dm}^2$}
{$2{,}6\text{ dm}^2$}
\loigiai{
Đổi $65\text{ cm} = 6{,}5\text{ dm}$.\\
Diện tích tấm tôn hình thoi là:
\[
\dfrac{8 \times 6{,}5}{2} = 26\text{ (dm}^2\text{)}.
\]
\textbf{Đáp án của câu hỏi là B}
}
\end{ex}

\begin{ex} %[5H3V16]
Cho bốn số đo diện tích: $A = 3\text{ m}^2\ 40\text{ dm}^2$; $B = 345\text{ dm}^2$; $C = 3\text{ m}^2\ 5\text{ dm}^2$; $D = 34\,000\text{ cm}^2$. Số đo lớn nhất là
\choice
{$A$}
{\True $B$}
{$C$}
{$D$}
\loigiai{
Đổi tất cả về cùng đơn vị đề-xi-mét vuông:
\begin{itemize}
    \item $A = 3\text{ m}^2\ 40\text{ dm}^2 = 340\text{ dm}^2$.
    \item $B = 345\text{ dm}^2$.
    \item $C = 3\text{ m}^2\ 5\text{ dm}^2 = 305\text{ dm}^2$.
    \item $D = 34\,000\text{ cm}^2 = 340\text{ dm}^2$.
\end{itemize}
So sánh: $305 < 340 = 340 < 345$. Số đo $B$ lớn nhất.\\
\textbf{Đáp án của câu hỏi là B}
}
\end{ex}

\begin{ex} %[5H3V16]
Một mặt bàn hình bình hành có độ dài đáy là $12\text{ dm}$, chiều cao tương ứng là $75\text{ cm}$. Diện tích mặt bàn đó tính theo mét vuông là
\choice
{$9\text{ m}^2$}
{\True $0{,}9\text{ m}^2$}
{$0{,}09\text{ m}^2$}
{$90\text{ m}^2$}
\loigiai{
Đổi sang đơn vị mét:
\begin{itemize}
    \item $12\text{ dm} = 1{,}2\text{ m}$.
    \item $75\text{ cm} = 0{,}75\text{ m}$.
\end{itemize}
Diện tích mặt bàn hình bình hành là:
\[
1{,}2 \times 0{,}75 = 0{,}9\text{ (m}^2\text{)}.
\]
\textbf{Đáp án của câu hỏi là B}
}
\end{ex}

\begin{ex} %[5H3V16]
Số thích hợp điền vào dấu hỏi chấm để $2\text{ dm}^2\ 4\text{ mm}^2 = ?\text{ mm}^2$ là
\choice
{$204$}
{$2\,004$}
{\True $20\,004$}
{$200\,004$}
\loigiai{
Vì $1\text{ dm}^2 = 100\text{ cm}^2 = 10\,000\text{ mm}^2$ nên:
\[
2\text{ dm}^2\ 4\text{ mm}^2 = 20\,000\text{ mm}^2 + 4\text{ mm}^2 = 20\,004\text{ mm}^2.
\]
\textbf{Đáp án của câu hỏi là C}
}
\end{ex}

\begin{ex} %[5H3V16]
Một tờ giấy màu hình chữ nhật có diện tích là $6\text{ dm}^2$. Bạn Lan cắt ra $8$ hình vuông nhỏ bằng nhau, mỗi hình vuông có diện tích $50\text{ cm}^2$. Diện tích phần giấy màu còn lại là
\choice
{$400\text{ cm}^2$}
{\True $200\text{ cm}^2$}
{$2\text{ cm}^2$}
{$20\text{ cm}^2$}
\loigiai{
Đổi $6\text{ dm}^2 = 600\text{ cm}^2$.\\
Tổng diện tích $8$ hình vuông nhỏ đã cắt là:
\[
8 \times 50 = 400\text{ (cm}^2\text{)}.
\]
Diện tích phần giấy màu còn lại là:
\[
600 - 400 = 200\text{ (cm}^2\text{)}.
\]
\textbf{Đáp án của câu hỏi là B}
}
\end{ex}

\begin{ex} %[5H3V16]
Người ta cấy lúa thí nghiệm trên một luống đất hình chữ nhật rộng $80\text{ cm}$, dài $25\text{ dm}$. Diện tích luống đất đó tính theo mét vuông là
\choice
{$20\text{ m}^2$}
{\True $2\text{ m}^2$}
{$0{,}2\text{ m}^2$}
{$200\text{ m}^2$}
\loigiai{
Đổi về mét:
\begin{itemize}
    \item $80\text{ cm} = 0{,}8\text{ m}$.
    \item $25\text{ dm} = 2{,}5\text{ m}$.
\end{itemize}
Diện tích luống đất là:
\[
0{,}8 \times 2{,}5 = 2\text{ (m}^2\text{)}.
\]
\textbf{Đáp án của câu hỏi là B}
}
\end{ex}

\begin{ex} %[5H3V16]
Một mảnh kính hình chữ nhật có chu vi là $28\text{ dm}$. Chiều dài hơn chiều rộng $4\text{ dm}$. Diện tích mảnh kính đó tính theo xăng-ti-mét vuông là
\choice
{$4\,500\text{ cm}^2$}
{\True $4\,500\text{ cm}^2$}
{$450\text{ cm}^2$}
{$45\text{ cm}^2$}
\loigiai{
Nửa chu vi mảnh kính là: $28 : 2 = 14\text{ (dm)}$.\\
Chiều dài mảnh kính: $(14 + 4) : 2 = 9\text{ (dm)}$.\\
Chiều rộng mảnh kính: $14 - 9 = 5\text{ (dm)}$.\\
Diện tích mảnh kính: $9 \times 5 = 45\text{ (dm}^2\text{)}$.\\
Đổi sang xăng-ti-mét vuông: $45\text{ dm}^2 = 4\,500\text{ cm}^2$.\\
\textbf{Đáp án của câu hỏi là A}
}
\end{ex}

\begin{ex} %[5H3V16]
Một hình thoi có diện tích bằng diện tích một hình vuông cạnh $6\text{ dm}$. Biết một đường chéo của hình thoi dài $9\text{ dm}$. Độ dài đường chéo kia của hình thoi là
\choice
{$4\text{ dm}$}
{\True $8\text{ dm}$}
{$6\text{ dm}$}
{$12\text{ dm}$}
\loigiai{
Diện tích hình vuông (và hình thoi) là:
\[
6 \times 6 = 36\text{ (dm}^2\text{)}.
\]
Độ dài đường chéo kia của hình thoi là:
\[
36 \times 2 : 9 = 8\text{ (dm)}.
\]
\textbf{Đáp án của câu hỏi là B}
}
\end{ex}

\textbf{Phần 2. Trắc nghiệm đúng sai}
\setcounterpageref{ex}{0}

\begin{ex} %[5H3V16]
Một căn phòng hình chữ nhật có kích thước $5\text{ m} \times 4\text{ m}$. Người ta dự định lát nền bằng hai phương án:
- Phương án 1: Dùng gạch men hình vuông cạnh $40\text{ cm}$
- Phương án 2: Dùng gạch men hình vuông cạnh $50\text{ cm}$
Xét tính đúng sai của các khẳng định sau:
\choiceTF
{\True Diện tích nền căn phòng là $20\text{ m}^2 = 200\,000\text{ cm}^2$.}
{\True Nếu dùng gạch cạnh $40\text{ cm}$ thì cần đúng $125$ viên gạch.}
{Nếu dùng gạch cạnh $50\text{ cm}$ thì cần $100$ viên gạch.}
{\True Số viên gạch cạnh $40\text{ cm}$ cần dùng nhiều hơn số viên gạch cạnh $50\text{ cm}$ là $45$ viên.}
\loigiai{
\begin{itemize}
    \item Diện tích nền: $5 \times 4 = 20\text{ m}^2 = 200\,000\text{ cm}^2$. Khẳng định a \textbf{Đúng}.
    \item Diện tích gạch $40\text{ cm}$: $40 \times 40 = 1\,600\text{ cm}^2$. Số viên cần: $200\,000 : 1\,600 = 125$ (viên). Khẳng định b \textbf{Đúng}.
    \item Diện tích gạch $50\text{ cm}$: $50 \times 50 = 2\,500\text{ cm}^2$. Số viên cần: $200\,000 : 2\,500 = 80$ viên $
e 100$ viên. Khẳng định c \textbf{Sai}.
    \item Chênh lệch số gạch: $125 - 80 = 45$ (viên). Khẳng định d \textbf{Đúng}.
\end{itemize}
\begin{itemize}
    \item a) \textbf{Đúng}.
    \item b) \textbf{Đúng}.
    \item c) \textbf{Sai}: Cần 80 viên.
    \item d) \textbf{Đúng}.
\end{itemize}
}
\end{ex}

\begin{ex} %[5H3V16]
Cho các số đo diện tích sau:
$M = 4\text{ m}^2\ 5\text{ dm}^2$ ; $N = 4{,}5\text{ m}^2$ ; $P = 405\text{ dm}^2$ ; $Q = 45\,000\text{ cm}^2$
Xét tính đúng sai của các khẳng định sau:
\choiceTF
{\True Đổi $M = 405\text{ dm}^2$.}
{\True Hai số đo $M$ và $P$ có giá trị bằng nhau.}
{\True Hai số đo $N$ và $Q$ có giá trị bằng nhau.}
{Số đo $M$ lớn hơn số đo $N$.}
\loigiai{
\begin{itemize}
    \item $M = 4\text{ m}^2\ 5\text{ dm}^2 = 405\text{ dm}^2$. Khẳng định a \textbf{Đúng}.
    \item $M = P = 405\text{ dm}^2$. Khẳng định b \textbf{Đúng}.
    \item $N = 4{,}5\text{ m}^2 = 450\text{ dm}^2 = 45\,000\text{ cm}^2 = Q$. Khẳng định c \textbf{Đúng}.
    \item Vì $405\text{ dm}^2 < 450\text{ dm}^2$ nên $M < N$. Do đó khẳng định $M > N$ là \textbf{Sai}.
\end{itemize}
\begin{itemize}
    \item a) \textbf{Đúng}.
    \item b) \textbf{Đúng}.
    \item c) \textbf{Đúng}.
    \item d) \textbf{Sai}.
\end{itemize}
}
\end{ex}

\textbf{Phần 3. Trả lời ngắn}
\setcounterpageref{ex}{0}

\begin{ex} %[5H3V16]
Đổi số đo $6\text{ m}^2\ 15\text{ dm}^2$ ra đơn vị đề-xi-mét vuông. Kết quả là bao nhiêu đề-xi-mét vuông?
\shortans{615}
\loigiai{
\textbf{Đáp số:} $615$\\
Vì $1\text{ m}^2 = 100\text{ dm}^2$ nên:
\[
6\text{ m}^2\ 15\text{ dm}^2 = 600\text{ dm}^2 + 15\text{ dm}^2 = 615\text{ dm}^2.
\]
}
\end{ex}

\begin{ex} %[5H3V16]
Một bảng thông báo hình chữ nhật có kích thước $120\text{ cm} \times 80\text{ cm}$. Diện tích của bảng thông báo đó bằng bao nhiêu đề-xi-mét vuông?
\shortans{96}
\loigiai{
\textbf{Đáp số:} $96$\\
Đổi sang đề-xi-mét:
\[
120\text{ cm} = 12\text{ dm}; \quad 80\text{ cm} = 8\text{ dm}.
\]
Diện tích bảng thông báo là:
\[
12 \times 8 = 96\text{ (dm}^2\text{)}.
\]
}
\end{ex}

\begin{ex} %[5H3V16]
Một bức tranh hình vuông có diện tích $6\,400\text{ cm}^2$. Hãy tính chu vi của bức tranh đó theo đơn vị đề-xi-mét.
\shortans{32}
\loigiai{
\textbf{Đáp số:} $32$\\
Đổi $6\,400\text{ cm}^2 = 64\text{ dm}^2$.\\
Vì $8 \times 8 = 64$ nên cạnh bức tranh hình vuông là $8\text{ dm}$.\\
Chu vi của bức tranh là:
\[
8 \times 4 = 32\text{ (dm)}.
\]
}
\end{ex}

\begin{ex} %[5H3V16]
Người ta lát sàn một phòng học diện tích $48\text{ m}^2$ bằng các tấm gỗ công nghiệp hình chữ nhật có diện tích $800\text{ cm}^2$ mỗi tấm. Hỏi cần dùng bao nhiêu tấm gỗ như thế?
\shortans{600}
\loigiai{
\textbf{Đáp số:} $600$\\
Đổi $48\text{ m}^2 = 480\,000\text{ cm}^2$.\\
Số tấm gỗ cần dùng là:
\[
480\,000 : 800 = 600\text{ (tấm)}.
\]
}
\end{ex}

\begin{ex} %[5H3V16]
Tìm số tự nhiên $x$ thỏa mãn:
\[
3\text{ m}^2\ 25\text{ dm}^2 < x\text{ dm}^2 < 3\text{ m}^2\ 27\text{ dm}^2.
\]
\shortans{326}
\loigiai{
\textbf{Đáp số:} $326$\\
Đổi sang đề-xi-mét vuông:
\[
325\text{ dm}^2 < x\text{ dm}^2 < 327\text{ dm}^2 \implies 325 < x < 327.
\]
Số tự nhiên duy nhất thỏa mãn là $x = 326$.
}
\end{ex}

\begin{ex} %[5H3V16]
Một hình thoi có diện tích $1\,500\text{ mm}^2$, độ dài một đường chéo là $5\text{ cm}$. Hãy tính độ dài đường chéo còn lại theo đơn vị mi-li-mét.
\shortans{60}
\loigiai{
\textbf{Đáp số:} $60$\\
Đổi $5\text{ cm} = 50\text{ mm}$.\\
Độ dài đường chéo còn lại là:
\[
1\,500 \times 2 : 50 = 60\text{ (mm)}.
\]
}
\end{ex}
